%% ch5_implementation.tex  —  Chapter 5: Implementation (rewrite staging)
%% HSP-Agile Thesis, East China Normal University, 2026
%% Ye Xujun

\chapter{Implementation}
\label{ch:implementation}

This chapter realises SpecIR and the SgDP loop as a deployable system:
which module owns which failure mode, what can be swapped for ablation, and
where latency is spent.
Algorithms and acceptance semantics remain in Chapter~\ref{ch:method};
here the concern is software realisation---adapters, checker wiring,
Semantic Feedback IR serialisation, and the conjunctive release gate.
Notation follows~\eqref{eq:accept}:
$t = (\mathcal{S}_t, \mathcal{X}_t, \mathcal{Y}_t)$, candidate~$c$,
$\mathrm{Conf}(c,t)$, and $\mathrm{Accept}(c,t)$.

The primary vignette is the tier classifier of Chapter~\ref{ch:intro}:
\texttt{OTHERWISE} must return \texttt{Review}, yet a smooth candidate
emits \texttt{Pass} while busy Reject/Pass tests stay green.
On that residual, the same modules build a Z3 catch-all witness, emit a
typed Semantic Feedback IR record, and decide $\mathrm{Accept}$.
Ordered-overlap tasks such as \code{classify\_signal} exercise the same
path as a secondary regime.

% ---------------------------------------------------------------------------
\section{System Architecture Overview}
\label{sec:impl:architecture}
% ---------------------------------------------------------------------------

Stages~1--4 must be ablatable and auditable without rewriting the rest
(Chapter~\ref{ch:experiments}), so ingestion, orchestration, and external
services stay apart.
Figure~\ref{fig:component-arch} shows the layout: pipeline core at the
centre, LLM/SMT clients below, and experiment orchestration above---the
software view of Figure~\ref{fig:sgdp-overview}.

\begin{figure}[t]
  \centering
  \IfFileExists{diagrams/puml/rendered/component_arch.pdf}{%
    \pumlfig{diagrams/puml/rendered/component_arch}%
  }{%
    \fbox{\parbox{0.90\linewidth}{\centering\vspace{1.2cm}
      \textit{[Figure pending render: component\_arch.pdf]}\\
      HSP-Agile system component architecture.
      \vspace{1.2cm}}}%
  }
  \caption{Component architecture: Stages~1--4 as isolatable modules under one
    orchestrator.
    Ablations swap a stage without rewriting the rest (software view of
    Figure~\ref{fig:sgdp-overview}).}
  \label{fig:component-arch}
\end{figure}

\texttt{PipelineRunner} owns the loop: an immutable \texttt{PipelineConfig},
delegation to \texttt{ECNUClient}, \texttt{FormalChecker}, and
\texttt{PatternMatcher} (Figure~\ref{fig:system-class}).
Value objects carry serialisable audit traces.
E8 adapters (Mini-Z, Mini-StateMachine) feed the same SpecIR path; typed
result objects keep stages isolatable.

\begin{figure}[t]
  \centering
  \IfFileExists{diagrams/puml/rendered/system_class.pdf}{%
    \pumlfig{diagrams/puml/rendered/system_class}%
  }{%
    \fbox{\parbox{0.90\linewidth}{\centering\vspace{1.2cm}
      \textit{[Figure pending render: system\_class.pdf]}\\
      HSP-Agile system class diagram.
      \vspace{1.2cm}}}%
  }
  \caption{Class diagram: one \texttt{run(task)} entry point composes
    Stages~1, 2, and~4 behind Algorithm~\ref{alg:sgdp}.
    \texttt{PipelineRunner} is the hub; clients and checkers are leaves.}
  \label{fig:system-class}
\end{figure}

\paragraph{Module responsibilities.}
Adapters keep notation changes out of the core: SOFL, Mini-Z, and
Mini-StateMachine each implement \texttt{SpecAdapter.parse() $\to$ SpecIR},
then lower via \texttt{FSFLowerer}.
The pipeline runner realises Algorithm~\ref{alg:sgdp}, including early exit on
$\mathrm{Accept}(c_k,t)=1$ and best-effort fallback when $k > K$.
The formal checker builds Z3 queries, executes candidates in a sandbox, and
assembles counterexamples.
The pattern guard applies catalogue $\mathcal{P}$ (Table~\ref{tab:patterns})
with AST and regex detectors.
The experiment orchestrator schedules mode--task pairs across a process pool
with resume-safe logs.

% ---------------------------------------------------------------------------
\section{Specification Ingestion Layer}
\label{sec:impl:ingestion}
% ---------------------------------------------------------------------------

The checker never consumes raw SOFL, Mini-Z, or Mini-StateMachine text.
Ingestion parses each surface notation into SpecIR
(Chapter~\ref{ch:formalization}) and lowers it to the FSF-shaped
\texttt{TaskSpec} consumed by Stages~2--4:
notation-specific parsing via a registered \texttt{SpecAdapter};
canonicalisation into ordered \texttt{GuardedCase} records; and lowering via
\texttt{FSFLowerer}.
The primary SOFL path uses the Agile-SOFL
toolchain~\citep{li2022qrs_companion,liu2004soflbook}; E8 adapters bypass SOFL
syntax while producing the same SpecIR structure.
Figure~\ref{fig:fsf-parser-class} shows the SOFL-side bridge classes.

\subsection{SpecIR Adapter Registry}
\label{sec:impl:specir-adapters}

The adapter registry (\path{src/adapters/__init__.py}) maps notation names to
concrete \texttt{SpecAdapter} subclasses.
Each adapter implements \texttt{parse(spec\_text, task\_id) $\to$ SpecIR},
populating \texttt{GuardedCase} objects with normalised \texttt{GuardAtom}
predicates, postconditions, and surface-text fields for prompting.
\texttt{SpecIR.to\_dict()} / \texttt{from\_dict()} round-trip through
\path{schemas/spec_ir.schema.json} for audit logging and cross-notation
interchange.

\texttt{FSFLowerer.lower(spec)} projects SpecIR onto the legacy
\texttt{TaskSpec} schema (ordered scenarios, I/O signatures, prompt text,
pattern-guard metadata).
\texttt{PipelineRunner.run()} accepts either a \texttt{TaskSpec} dict or a
\texttt{SpecIR} instance (\path{src/pipeline/task_input.py}), normalising
through the lowerer when needed.
Oracle construction, witness generation, and repair feedback therefore stay
notation-agnostic while remaining compatible with the 120-task FSF benchmark.

\begin{figure}[t]
  \centering
  \IfFileExists{diagrams/puml/rendered/fsf_parser_class.pdf}{%
    \pumlfig{diagrams/puml/rendered/fsf_parser_class}%
  }{%
    \fbox{\parbox{0.90\linewidth}{\centering\vspace{1.2cm}
      \textit{[Figure pending render: fsf\_parser\_class.pdf]}\\
      FSF specification parsing and test-case generation class diagram.
      \vspace{1.2cm}}}%
  }
  \caption{Surface syntax stops at SpecIR; Stages~2--4 consume only the
           lowered \texttt{TaskSpec}.
           For the tier classifier, the bridge yields a three-scenario task
           whose catch-all residual is $S_3$ (\texttt{OTHERWISE}).}
  \label{fig:fsf-parser-class}
\end{figure}

\paragraph{Parsing sketch.}
The Agile-SOFL bridge extracts, per process with a non-empty FSF block, a
\texttt{TaskSpec}: task id, callable signature, ordered
\texttt{FSFScenario} list (guard \texttt{test}, assignment \texttt{def},
kind \texttt{scenario}/\texttt{others}), external-variable map, and rendered
prompt text.
Scenario indices preserve first-match order
$s_1 \succ \cdots \succ s_n$~\citep{dijkstra1975guarded}
(Definition~\ref{def:scenario}).
Guards are normalised to a shared relational-atom syntax consumed by the
prompt renderer, Z3 translator, and pattern guard; output names define
$\mathcal{Y}_t$ for field-wise oracle comparison and RF04 detection.
External \texttt{ext rd}/\texttt{ext wr} names are listed in the prompt and
checked by RF10 when relevant.

% ---------------------------------------------------------------------------
\section{LLM Integration Layer}
\label{sec:impl:llm}
% ---------------------------------------------------------------------------

Synthesis is the only non-deterministic stage; everything else must be
replayable.
The LLM layer wraps an OpenAI-compatible gateway
(\texttt{ECNUClient}, \path{src/llm/ecnu_client.py}) for auth, retry, timeout,
token accounting, and audit logs---the name is historical, not a coupling.
Production runs use Qwen2.5-27B-Instruct (\textsf{Qwen-27B}) with exponential
backoff (up to three retries), a 90\,s request timeout, and optional JSON
audit of message history and token counts~\citep{openai2023gpt4,chen2021codex}.

Prompt assembly follows Section~\ref{subsec:method:prompt}: a fixed system
message (role, fenced-code output, first-match guard order) plus a user
message carrying the rendered FSF text, signature, required output keys, and---on
repair iterations---verbalised feedback.
Sampling uses $T_\text{gen}{=}0.7$ on the first attempt and
$T_\text{repair}{=}0.0$ thereafter ($\mathit{top\_p}{=}0.95$): explore once,
then exploit under a deterministic checker.
Responses are scanned for the last fenced Python block (or a conservative
line heuristic); \texttt{compile()} failure yields a sentinel~$\bot$ that
scores zero and triggers repair without crashing the pipeline.

\subsection{Semantic Feedback IR and Two-Layer Repair Prompting}
\label{sec:impl:repair-feedback}

When iteration $k$ fails formal checking or pattern screening, the repair
engine constructs a \texttt{SemanticFeedbackIR}
(\path{src/repair/feedback_ir.py}) from checker counterexamples and the
active \texttt{TaskSpec}.
Each record serialises to JSON per
\texttt{schemas/semantic\_feedback.schema.json} \emph{before} any
natural-language rendering:
\begin{enumerate}
  \item \textbf{IR construction} --- nine typed fields
        (\texttt{violation\_type}, \texttt{scenario\_index},
        \texttt{constraint\_text}, \texttt{inputs}, \texttt{expected},
        \texttt{observed}, \texttt{reason}, \texttt{priority},
        \texttt{suggested\_fix}) from witness and scenario metadata;
  \item \textbf{JSON serialisation} --- schema-compliant records logged in
        \texttt{PipelineResult.semantic\_feedback};
  \item \textbf{Prompt projection} --- \texttt{FeedbackRenderer.render()}
        maps the same IR to \texttt{test\_only} (B2 / E6-A),
        \texttt{test\_expected} (E6-B), or \texttt{semantic\_ir} (M / E6-C).
\end{enumerate}

Pattern violations append as a capped list of \texttt{pattern\_id:name}
pairs under a \emph{``Previous attempt failed''} header.
Ablation E6 varies only the renderer while holding IR construction fixed,
so Semantic Feedback is a genuine intermediate representation rather than a
prompt template alone.

\paragraph{Example: catch-all failure.}
On the tier classifier (Listings~\ref{lst:others-fsf},~\ref{lst:others-bug}),
a mid-band witness activates $S_3$ (\texttt{OTHERWISE}).
The oracle requires \texttt{Review}; the candidate returns \texttt{Pass}.
The checker emits a record of the following shape (abbreviated):
\begin{lstlisting}[language={},basicstyle=\ttfamily\footnotesize,frame=single]
{
  "violation_type": "output",
  "scenario_index": 3,
  "constraint_text": "OTHERWISE",
  "inputs": {"score": 3, "floor": 10},
  "expected": {"tier": "Review"},
  "observed": {"tier": "Pass"},
  "reason": "catch-all residual returns wrong constant",
  "priority": 1,
  "suggested_fix": "return Review on OTHERWISE"
}
\end{lstlisting}
A \texttt{test\_only} projection would state only that the case failed;
\texttt{semantic\_ir} names the residual scenario, expected constant, and
repair hint---the implementation hinge behind the C2 mechanism in
Chapter~\ref{ch:method}.

% ---------------------------------------------------------------------------
\section{Formal Conformance Checker}
\label{sec:impl:checker}
% ---------------------------------------------------------------------------

Acceptance must not rest on lucky unit tests.
The formal checker is the deterministic oracle that grounds release decisions
in SMT-backed witnesses rather than heuristic
sampling~\citep{claessen2000quickcheck}.
Given candidate~$c$ and task~$t$, it computes $\mathrm{Conf}(c,t)$ and
assembles $\mathit{cexs}$ for repair (Figure~\ref{fig:checker-activity};
Algorithm~\ref{alg:checker} in Chapter~\ref{ch:method}).

\begin{figure}[t]
  \centering
  \IfFileExists{diagrams/puml/rendered/checker_activity.pdf}{%
    \pumlfig{diagrams/puml/rendered/checker_activity}%
  }{%
    \fbox{\parbox{0.90\linewidth}{\centering\vspace{1.2cm}
      \textit{[Figure pending render: checker\_activity.pdf]}\\
      Formal conformance checker activity diagram.
      \vspace{1.2cm}}}%
  }
  \caption{Checker activity: case generation and evaluation are separate
           partitions; counterexamples feed Stage~3 repair.}
  \label{fig:checker-activity}
\end{figure}

Guards lower to Z3 integer constraints under the priority encoding
\eqref{eq:priority}:
$\Phi_i \equiv \widehat{g}_i \land \bigwedge_{j<i}\lnot\widehat{g}_j$,
with \texttt{others} as the negation of preceding guards
(Definition~\ref{def:oracle}).
Each $\Phi_i$ uses a fresh solver.
Integer variables are bounded to $[-5,20]$ for Z3 termination in the checker
implementation; hard-task metadata describes inputs over $[0,100]^5$
(Chapter~\ref{ch:experiments}), so reported witnesses use the implemented
solver bounds rather than the full declared domain
(Appendix~\ref{app:reproducibility}).
Satisfying models become concrete witnesses $(x,y,i)$; model blocking yields
up to a few alternatives per scenario (Algorithm~\ref{alg:case-gen}).
On the tier classifier, $\Phi_3$ forces a residual mid-band point that
Reject/Pass unit tests never need to hit.

Candidates run in an isolated namespace; field-wise integer equality against
the oracle yields
$\mathrm{Conf}(c,t) = |\mathit{passed}| / |\mathcal{D}(t)|$
\eqref{eq:conformance}.
Repair iterations use a 16-case budget ($N_{\max}$ for mode~M); post-pipeline
strict scoring uses 64 cases (\texttt{PipelineConfig}; sensitivity in
Section~\ref{sec:results:sensitivity}).

\begin{algorithm}[t]
\caption{Z3-Based Concrete Case Generation (sketch)}
\label{alg:case-gen}
\begin{algorithmic}[1]
\Require ordered scenarios $\mathcal{S}_t$, variables $V$, budget $N$
\Ensure case set $\mathcal{D}(t)$ with $|\mathcal{D}(t)| \le N$
\For{$i = 1$ \textbf{to} $n$}
  \State build $\Phi_i$ per~\eqref{eq:priority}; solve under domain bounds
  \While{$\mathit{sat}$ \textbf{and} budget remains}
    \State extract model $x$; set $y \leftarrow \phi_i(x)$;
           record $(x,y,i)$; block model and re-query
  \EndWhile
\EndFor
\State \Return $\mathcal{D}(t)$
\end{algorithmic}
\end{algorithm}

% ---------------------------------------------------------------------------
\section{Counterexample-Guided Repair Engine}
\label{sec:impl:repair}
% ---------------------------------------------------------------------------

The repair engine adapts CEGIS~\citep{solarlezama2008sketching} to an LLM
synthesiser: it keeps the feedback buffer, increments the attempt counter,
and exits on $\mathrm{Accept}$ or best-effort when $k > K$.
Control states are in Figure~\ref{fig:pipeline-state}; the feedback payload
on each repair arc is Figure~\ref{fig:repair-feedback}
(Chapter~\ref{ch:method}).

\begin{figure}[t]
  \centering
  \IfFileExists{diagrams/puml/rendered/pipeline_state.pdf}{%
    \pumlfig{diagrams/puml/rendered/pipeline_state}%
  }{%
    \fbox{\parbox{0.90\linewidth}{\centering\vspace{1.2cm}
      \textit{[Figure pending render: pipeline\_state.pdf]}\\
      HSP-Agile pipeline state machine.
      \vspace{1.2cm}}}%
  }
  \caption{Pipeline state machine: repair aggregates counterexamples and
           pattern hits into Semantic Feedback IR before re-synthesis.
           Acceptance requires formal conformance and a clear high-severity
           pattern screen.}
  \label{fig:pipeline-state}
\end{figure}

States are \textsc{Init}, \textsc{Synthesizing}, \textsc{FormalChecking},
\textsc{PatternChecking}, \textsc{Repairing}, and terminal
\textsc{Accepted}/\textsc{Exhausted}.
Formal failure or any high-severity pattern hit routes to repair; both
$\mathrm{Conf}=1$ and a clear $\mathcal{P}^H$ screen yield acceptance.
Budget $K{=}3$ (principle~P2) caps LLM calls; mode~A3 disables repair for a
single-shot path.
On exhaustion the pipeline returns the highest-$\mathrm{Conf}$ candidate as
\textsc{Best-Effort}.
Algorithm~\ref{alg:repair-loop} summarises the control logic; the full SgDP
loop is Algorithm~\ref{alg:sgdp}.

\begin{algorithm}[t]
\caption{Counterexample-Guided Repair Loop (sketch)}
\label{alg:repair-loop}
\begin{algorithmic}[1]
\Require task $t$, budget $K$, flags
         $(\mathit{enable\_formal}, \mathit{enable\_patterns}, \mathit{enable\_repair})$
\Ensure $c^*$, status $\in \{\textsc{Accept}, \textsc{Best-Effort}\}$
\State $\mathit{feedback} \leftarrow \emptyset$;\quad
       $\mathit{best} \leftarrow (\bot, 0)$
\For{$k = 1$ \textbf{to} $K$}
  \State $c_k \leftarrow \ProcCall{LLM}{t,\, \mathit{feedback}}$
  \State evaluate $\mathrm{Conf}$ / $\mathit{cexs}$ and $\mathit{hits}$
        under the enabled flags
  \State update $\mathit{best}$; \textbf{if}
        $\mathrm{Accept}(c_k,t)=1$ \textbf{then return} Accept
  \State \textbf{if} repair enabled \textbf{then}
        $\mathit{feedback} \leftarrow
        \ProcCall{BuildFeedback}{\mathit{cexs}_k,\, \mathit{hits}_k,\, k}$
\EndFor
\State \Return $\mathit{best}$, Best-Effort
\end{algorithmic}
\end{algorithm}

% ---------------------------------------------------------------------------
\section{Structural Screen Module}
\label{sec:impl:patterns}
% ---------------------------------------------------------------------------

A candidate can match every witness in $\mathcal{D}(t)$ and still hard-code
success, ignore ordered guards, or drop an output field that the finite case
set never exercised.
Stage~4 therefore implements the structural half of acceptance safety: an
oracle $\mathrm{Screen}(c,t)$ independent of the witness suite that can veto
release even when $\mathrm{Conf}{=}1$
(Section~\ref{sec:method:guard}; Figure~\ref{fig:accept-tree}).

The present artifact instantiates $\mathrm{Screen}$ with a catalogue of
requirement-level code smells $\mathcal{P}$ (RF01--RF14;
Table~\ref{tab:patterns}), in the spirit of FindBugs-style structural
screening~\citep{hovemeyer2004findbugs,ayewah2007findbugs} and prior SOFL
fault-prevention patterns~\citep{li2023iet_faultprevention}.
Detectors are primarily AST predicates (control flow, return shape, missing
fields); a minority use surface matching or guard--conditional
cross-reference.
The blocking high-severity subset is
\[
  \mathcal{P}^H = \{\text{RF01}, \text{RF02}, \text{RF04}, \text{RF05},
    \text{RF06}, \text{RF07}, \text{RF09}, \text{RF11}, \text{RF13}\};
\]
medium/low hits are logged into repair feedback but do not alone withhold
$\mathrm{Accept}$.
AST detectors catch structural shortcuts (e.g.\ RF04 missing outputs, RF07
unconditional success, RF09 branch-count mismatch); cross-reference
detectors catch guard-literal mismatches (RF01--RF02, RF05--RF06, RF11).
Each $\delta_p$ is a pure predicate; hits carry id, severity, location, and
message for the Semantic Feedback IR renderer.

The module boundary is intentionally thin: any object that maps $(c,t)$ to
severity-tagged hits can replace $\{\delta_p\}$ without touching SpecIR, Z3
witnesses, or Equation~\eqref{eq:accept}.
In a production Agile-SOFL pipeline the same slot can hold a project-specific
smell pack, an existing CI static analyser, or an LLM-driven review agent that
returns structured vetoes into the same repair channel.
RF01--RF14 remain the evaluated default so C3 (PDR/FAR) is measured under a
deterministic screen.

The structural screen is the second conjunct of $\mathrm{Accept}(c,t)$:
\[
  \mathrm{Accept}(c, t)
  \;=\;
  \mathbf{1}\!\left[\mathrm{Conf}(c, t) = 1\right]
  \;\land\;
  \mathbf{1}\!\left[\mathrm{Screen}(c, t) = \mathit{pass}\right].
\]
With the PoC encoding,
$\mathrm{Screen}(c,t)=\mathit{pass}$ iff no $p\in\mathcal{P}^H$ fires.
The screen runs \emph{after} formal checking on each iteration.
A candidate that passes all Z3 cases but triggers RF07 (unconditional
success) is routed to repair, with the smell verbalised alongside any
counterexamples.
Witnesses catch catch-all and ordering residuals; the screen catches
shortcuts that never attempt the required structure
(Section~\ref{subsec:conformance-acceptance}).

% ---------------------------------------------------------------------------
\section{Experiment Orchestration}
\label{sec:impl:orchestration}
% ---------------------------------------------------------------------------

The orchestrator prioritises parallel throughput, crash-safe resumption, and
live progress~\citep{harman2004maff,beck2004xp}.
Two matrices are used: the \emph{core} seven-mode matrix is
$7 \times 120 = 840$ runs (B0--B2, M, A1--A3 under
\code{run\_hard\_full\_parallel\_v1}); the \emph{extended} matrix is
$10 \times 120 = 1{,}200$ when B3--B5 are included
(\code{run\_ccf\_b\_extended\_v1}; B6 is reported separately where evaluated).
Jobs run in a \emph{process pool} (one worker per core) because Z3's native
global state is not thread-safe; each worker owns its own client, config, and
pipeline, then appends a strict 64-case evaluation (and optional mutation
score~\citep{demillo1978mutation,li2023iet_faultprevention}) to a shared JSONL log.
Incremental persistence plus a \texttt{progress.json} sidecar make mid-batch
restarts resume-safe without re-running finished keys.

Table~\ref{tab:impl-modes} lists the seven core modes; flags on
\texttt{PipelineConfig} enable ablations A1--A3 and baselines B0--B2 against
full mode~M (Chapter~\ref{ch:experiments}).

\begin{table}[H]
  \caption{Execution mode configuration matrix.  Flags control which
           pipeline stages are active; numeric parameters tune per-stage
           behaviour.  All LLM-based modes use Qwen2.5-27B-Instruct (\textsf{Qwen-27B})
           with $K = 3$ unless noted.
           $N_{\max}$ is the per-scenario formal case cap.}
  \label{tab:impl-modes}
  \centering
  \thesistable{%
    \footnotesize
    \renewcommand{\arraystretch}{1.15}
    \begin{tabular*}{\linewidth}{@{\extracolsep{\fill}}l l c c c r@{}}
      \toprule
      Mode & Description & Formal & Guard & Repair & $N_{\max}$ \\
      \midrule
      B0 & Reference template  & \textit{eval only} & No  & No  & --- \\
      B1 & One-shot LLM        & No  & No  & No  & --- \\
      B2 & Test-feedback loop  & No  & No  & Yes & 8 \\
      M  & Full HSP-Agile      & Yes & Yes & Yes & 16 \\
      A1 & No formal check     & No  & Yes & Yes & 10 \\
      A2 & No pattern guard    & Yes & No  & Yes & 24 \\
      A3 & No repair loop      & Yes & Yes & No  & 24 \\
      \bottomrule
    \end{tabular*}%
  }
  \par\smallskip
  \raggedright\footnotesize All modes apply strict evaluation with a
  64-case budget after pipeline completion.
  B0 uses a deterministic template renderer; no LLM calls are made.
\end{table}

% ---------------------------------------------------------------------------
\section{Design Decisions and Trade-offs}
\label{sec:impl:tradeoffs}
% ---------------------------------------------------------------------------

Three decisions dominate the rest~\citep{woodcock2009formalsurvey}.

\subsection{Z3 Over Random Testing}

Z3~\citep{demoura2008z3} is preferred to random or property-based
testing~\citep{claessen2000quickcheck,ammann2016testing} because
(i)~guard regions under conjunctions of inequalities are rarely hit by
uniform sampling over $\mathbb{Z}^5$;
(ii)~priority encoding $\Phi_i$ yields exclusive first-match witnesses,
including catch-all residuals that busy-band samples miss; and
(iii)~failures arrive as scenario-indexed concrete counterexamples usable by
CEGIS repair~\citep{sun2024clover}.
The cost is solver latency ($\sim$50--200\,ms per query; $\sim$2\,s per
16-case repair iteration; $\sim$6--8\,s for 64-case strict scoring).

\subsection{Conjunctive Acceptance Over Single-Metric Gating}

A disjunctive gate (accept on $\mathrm{Conf}=1$ \emph{or} a clear structural
screen) would maximise throughput but conflate distinct fault classes.
Formal conformance certifies pointwise witness correctness; $\mathrm{Screen}$
certifies that the candidate is not an obvious structural shortcut
(Section~\ref{sec:impl:patterns}).
Neither implies the other: hard-coded outputs can pass sparse witnesses while
failing RF07, and a structurally clean candidate can still fail on a catch-all
or overlap residual.
The conjunctive $\mathrm{Accept}$ of~\eqref{eq:accept} is therefore
conservative---fewer accepts, but each satisfies both criteria.
Which concrete $\mathrm{Screen}$ fills the slot does not change that design.

\subsection{Latency vs.\ Quality Trade-off}

The 16/64 case split and $K{=}3$ bound per-task cost for sprint-cycle use
(P2).
Sixteen internal cases give at least two witnesses per scenario on
eight-scenario tasks; sixty-four strict cases probe residual and overlap
regions more densely.
Sensitivity analysis (Chapter~\ref{ch:experiments}) shows diminishing
conformance returns beyond 32 internal cases; most repairable faults resolve
within two iterations, so larger $K$ rarely pays off.

% ---------------------------------------------------------------------------
\section*{Chapter Summary}
\addcontentsline{toc}{section}{Chapter Summary}
% ---------------------------------------------------------------------------

The realisation maps cleanly onto the Tier-1 claims:
\begin{itemize}
  \item \textbf{C1} --- SpecIR adapter registry and FSF lowering keep surface
        syntax out of the checker.
  \item \textbf{C2} --- Semantic Feedback IR and the repair state machine own
        typed failure state across iterations (catch-all residual as primary
        example).
  \item \textbf{C3} --- Z3 conformance, pluggable $\mathrm{Screen}$
        (default: $\mathcal{P}^H$ smells), and conjunctive $\mathrm{Accept}$
        decide release.
\end{itemize}
Orchestration and sampling parameters are supporting engineering.
Whether those choices pay off is empirical---Chapter~\ref{ch:experiments}.
